% Generated draft: requires mathematical and source review.
\documentclass[11pt]{article}
\usepackage[margin=25mm]{geometry}
\usepackage{fontspec}
\setmainfont{Latin Modern Roman}
\usepackage{amsmath,amssymb,mathtools,unicode-math}
\setmathfont{Latin Modern Math}
\usepackage{enumitem,needspace,xurl,hyperref}
\setlength{\parindent}{0pt}
\setlength{\parskip}{5pt}
\setlength{\emergencystretch}{3em}
\newcommand{\N}{\mathbb{N}}
\newcommand{\Z}{\mathbb{Z}}
\newcommand{\Q}{\mathbb{Q}}
\newcommand{\R}{\mathbb{R}}
\newcommand{\C}{\mathbb{C}}
\newcommand{\F}{\mathbb{F}}
\newcommand{\A}{\mathbb{A}}
\newcommand{\PP}{\mathbb{P}}
\newcommand{\E}{\mathbb{E}}
\newcommand{\eps}{\varepsilon}
\newcommand{\dd}{\,\mathrm{d}}
\begin{document}
\section*{UnsolvedMath ID: 1207}
\textbf{Source problem code:} catalogue entry
\subsection{Definitions and mathematical statement}
For an immersed hypersurface \(f:M^n\to S^{n+1}(1)\), the target is the unit round sphere. In a local choice of unit normal, let \(A\) be the shape operator with principal curvatures \(\kappa_1,\ldots,\kappa_n\). Set
\[
H=\frac1n\operatorname{tr}A=\frac1n\sum_{i=1}^n\kappa_i,
\qquad S=|A|^2=\sum_{i=1}^n\kappa_i^2.
\]
Minimality means that the mean-curvature vector vanishes, equivalently \(H=0\) in every local normal choice. Thus the literal question, ``What minimal hypersurfaces in spheres have constant mean curvature?'', imposes no additional condition on a minimal hypersurface. It does not determine a nontrivial classification conjecture.

The supplied triage dated 2026-08-17 distinguishes the intended scalar-curvature question. For a minimal hypersurface in the unit sphere, the scalar curvature \(R\) satisfies
\[
R=n(n-1)-S.
\]
Hence constant scalar curvature is equivalent to constant squared norm of the second fundamental form. A hypersurface is isoparametric when its principal curvatures are constant, with the usual local-normal interpretation when necessary. The stronger standard Chern conjecture described by the supplied triage asks whether a closed minimal hypersurface in a sphere with constant scalar curvature must be isoparametric. Here closed means compact without boundary. The source has not fixed whether the intended category is immersed or embedded, nor the dimension range or connectedness conventions; these must be fixed before that alternative is an exact replacement statement. The source also mentions a distinct weaker question about discreteness of the possible constant scalar-curvature values, and does not choose between that question and isoparametric classification.

The background's suggestion that the totally geodesic sphere is the only minimal constant-scalar-curvature example is explicitly rejected by the supplied triage, which points to Clifford and other isoparametric examples. The triage reports progress in specified proper-Dupin cases due to Miyaoka and results in \(S^5(1)\) under additional constant Gauss--Kronecker curvature or related hypotheses. Gauss--Kronecker curvature is the product of the principal curvatures. These restricted results are not stated to resolve the general conjecture. The registry status is open; the supplied dated triage labels the conflated row partially solved and the intended general scalar-curvature conjecture open.
\subsection{Short English statement}
Every minimal hypersurface already has constant mean curvature zero, so the written question adds no condition. The supplied notes instead discuss Chern's conjectures about constant scalar curvature, discreteness of curvature values, and isoparametricity; the intended invariant and hypersurface category must be fixed.
\subsection{Sources}
\begin{itemize}
\item Local UnsolvedMath registry, permanent ID 1207, source code catalogue entry, ``Chern's Conjecture for Hypersurfaces in Spheres,'' including background attribution to Shiing-Shen Chern and supplied literature triage dated 2026-08-17.
\item Reiko Miyaoka, ``Chern's Conjecture in the Dupin case,'' arXiv:2504.02621 (2025): \url{https://arxiv.org/abs/2504.02621}. The supplied triage cites the scalar-curvature/isoparametric formulation and restricted proper-Dupin results.
\item Qintao Deng and Yunjia Kou, ``Closed minimal hypersurfaces in \(S^5(1)\) with constant scalar and Gauss-Kronecker curvatures,'' arXiv:2607.06588 (2026): \url{https://arxiv.org/abs/2607.06588}. The supplied triage cites classification under the additional constant Gauss--Kronecker hypothesis.
\end{itemize}
\end{document}
